(b) 确定幂变换 $\hat{\lambda}_{1}$ 使变量 $x_{1}$ 接近正态. 画出变换后观测值的 Q-Q 图.

(c) 确定幂变换 $\hat{\lambda}_{2}$ 使变量 $x_{2}$ 接近正态. 画出变换后观测值的 Q-Q 图.

(d) 用\hyperref[eq:4-40]{式(4-40)} 确定接近二元正态的幂变换.

4.41\quad 考虑练习 3.20 中的扫雪数据.

(a) 对原始数据的散点图中所有可能的离群值作出评论.

(b) 确定幂变换 $\hat{\lambda}_{1}$ 使变量 $x_{1}$ 接近正态. 画出变换后观测值的 Q-Q 图.

(c) 确定幂变换 $\hat{\lambda}_{2}$ 使变量 $x_{2}$ 接近正态. 画出变换后观测值的 Q-Q 图.

(d) 用\hyperref[eq:4-40]{式(4-40)} 确定接近二元正态的幂变换.

\phantomsection\label{bib:4}
\section*{参考文献}

\begin{enumerate}[leftmargin=2em,itemsep=2pt]
\item Anderson, T. W. \textit{An Introduction to Multivariate Statistical Analysis} (3rd ed.). New York: John Wiley, 2003.
\item Andrews, D. F., R. Gnanadesikan, and J. L. Warner. “Transformations of Multivariate Data.” \textit{Biometrics}, \textbf{27}, no. 4 (1971), 825--840.
\item Box, G. E. P., and D. R. Cox. “An Analysis of Transformations” (with discussion). \textit{Journal of the Royal Statistical Society (B)}, \textbf{26}, no. 2 (1964), 211--252.
\item Daniel, C. and F. S. Wood, \textit{Fitting Equations to Data: Computer Analysis of Multifactor Data}. New York: John Wiley, 1980.
\item Filliben, J. J. “The Probability Plot Correlation Coefficient Test for Normality.” \textit{Technometrics}, \textbf{17}, no. 1 (1975), 111--117.
\item Gnanadesikan, R. \textit{Methods for Statistical Data Analysis of Multivariate Observations} (2nd ed.). New York: Wiley-Interscience, 1977.
\item Hawkins, D. M. \textit{Identification of Outliers}. London, UK: Chapman and Hall, 1980.
\item Hernandez, F., and R. A. Johnson. “The Large-Sample Behavior of Transformations to Normality.” \textit{Journal of the American Statistical Association}, \textbf{75}, no. 372 (1980), 855--861.
\item Hogg, R. V., Craig, A. T. and J. W. Mckean \textit{Introduction to Mathematical Statistics} (6th ed.). Upper Saddle River, N.J.: Prentice Hall, 2004.
\item Looney, S. W., and T. R. Gulledge, Jr. “Use of the Correlation Coefficient with Normal Probability Plots.” \textit{The American Statistician}, \textbf{39}, no. 1 (1985), 75--79.
\item Mardia, K. V., Kent, J. T. and J. M. Bibby. \textit{Multivariate Analysis} (Paperback). London: Academic Press, 2003.
\item Shapiro, S. S., and M. B. Wilk. “An Analysis of Variance Test for Normality (Complete Samples).” \textit{Biometrika}, \textbf{52}, no. 4 (1965), 591--611.
\item Verrill, S., and R. A. Johnson. “Tables and Large-Sample Distribution Theory for Censored-Data Correlation Statistics for Testing Normality.” \textit{Journal of the American Statistical Association}, \textbf{83}, no. 404 (1988), 1192--1197.
\item Yeo, I. and R. A. Johnson “A New Family of Power Transformations to Improve Normality or Symmetry.” \textit{Biometrika}, \textbf{87}, no. 4 (2000), 954--959.
\item Zehna, P. “Invariance of Maximum Likelihood Estimators.” \textit{Annals of Mathematical Statistics}, \textbf{37}, no. 3 (1966), 744.
\end{enumerate}
